\documentclass[10pt,letterpaper]{article}
\usepackage[T1]{fontenc}
\usepackage{lmodern}
\usepackage[margin=25mm]{geometry}
\usepackage{amsmath,amssymb,amsthm,booktabs,array,tabularx}
\usepackage[hidelinks]{hyperref}
\usepackage{xurl}
\newtheorem{theorem}{Theorem}
\newtheorem{proposition}[theorem]{Proposition}
\newtheorem{lemma}[theorem]{Lemma}
\renewcommand{\thetheorem}{S\arabic{theorem}}
\renewcommand{\thesection}{S\arabic{section}}
\setlength{\emergencystretch}{3em}
\newcommand{\E}{\mathbb E}
\newcommand{\ind}{\mathbf 1}
\newcommand{\PASS}{\mathsf{PASS}}
\newcommand{\UNKNOWN}{\mathsf{UNKNOWN}}
\title{Technical Supplement to\\Adaptive HACT: Information Limits and Safe Adaptation of Verification Evidence}
\author{Da Xu and Fei Wang\\China Mobile Research Institute\\
\small Corresponding author: \href{mailto:xudayj@chinamobile.com}{xudayj@chinamobile.com}}
\date{}
\hypersetup{pdftitle={Adaptive HACT: Technical Supplement},pdfauthor={Da Xu and Fei Wang}}
\begin{document}
\maketitle
This supplement accompanies the primary manuscript. It supplies the full
complete-epoch orbit certificate, fixed-width padding proof and a practical
route through the reproducibility artifact. It is supporting material for
one submission. All numerical observations retain their original provenance;
no new checker campaign, consumer study or agent evaluation is introduced.

\section{Notation and scope}
A rooted tree represents a registry of $n$ distinct checks. Internal arities
lie in $\{2,\ldots,b\}$, and a packet costs $c(d)=H+Sd$ bytes.
A node with semantic coverage $G$ is emitted once for each completion wave
that intersects $G$. A strict epoch additionally constructs every internal
packet once. Thus $W(T)=\sum_v c(d_v)$ is construction cost and
$h_e(G)$ counts touched waves. Acceptance still requires every current
approved check to pass; none of these costs changes that predicate.
The main article gives the parameter-general persistent-refresh separation
and the global exchangeable optimum. The following two proofs complete its
complete-epoch and fixed-width extensions.

\section{Complete-epoch separation and orbit certificate}
\subsection{The separation also occurs in complete epochs}
The preceding construction uses persistent refresh events. A complete
verification epoch imposes more structure: the completion sets are disjoint
and cover every declared check. That restriction does not make pairwise
summaries sufficient.

\begin{theorem}[Complete-epoch information separation]\label{thm:strict-pairwise}
For any integer fanout $b\ge6$, $S>0$ and
$16/11<H/S<28/17$, there are two laws on six checks partitioned into exactly
three nonempty refresh windows, with one full initial construction, that
have identical population singleton and pairwise completion statistics
but disjoint unique global tree optima.
Writing
\[
 \Delta_P=(28S-17H)/8,\qquad
 \Delta_Q=(11H-16S)/5,
\]
the sharp deterministic pairwise-only regret is
$\min(\Delta_P,\Delta_Q)$; the sharp randomized minimax regret is
$\Delta_P\Delta_Q/(\Delta_P+\Delta_Q)$.
\end{theorem}
\begin{proof}
Let $G=\{0,1,2,3\}$ and $O=\{4,5\}$.
A block's type $(g,o)$ counts its members from $G$ and $O$.
Let $A=(G,\{4\},\{5\})$, and let $B,C,D$ be uniform laws over all
partitions of the respective block-type multisets
\begin{align*}
 B&:\{(2,1),(2,0),(0,1)\},\\
 C&:\{(2,1),(1,1),(1,0)\},\\
 D&:\{(2,0),(1,1),(1,1)\}.
\end{align*}
There are 12, 24 and 12 partitions in those families. Define
\[
 P=\tfrac78 A+\tfrac18 B,\qquad
 Q=\tfrac9{10}A+\tfrac1{20}C+\tfrac1{20}D.
\]
Every check completes once in every outcome. Each within-$G$ pair shares
a window with probability $11/12$, each $G$--$O$ pair with probability
$1/32$, and the $O$ pair with probability zero, under either law.
Uniformly permuting each partition's three windows also makes every
window-specific first and second moment agree. The expected number of
windows touching $G$ is nevertheless $9/8$ under $P$ and $6/5$ under $Q$.

We next reduce \emph{all} tree shapes and leaf permutations.
A node's weight, including its initial packet, lies in $[2,4]$.
Contracting a nonroot internal child of arity $d=2$ or $3$ raises the
parent's price by $S(d-1)$ per emission and removes that child's packet.
Its cost change in every complete three-window outcome is at most
\[
 4S(d-1)-2(H+Sd),
\]
which is $-2H$ for $d=2$ and $-2(H-S)$ for $d=3$, both negative.
Contraction stays feasible because there are only six leaves and $b\ge6$.
Since $\sum_v(d_v-1)=5$, a tree with at least three internal nodes
must contain such a nonroot child. Repeated contractions leave a flat root
or one root with a single internal child. Arity-two/three children are
themselves dominated by the flat root.

The remaining five child types are listed in
Table~\ref{tab:strict-orbits}, with $x=H/S$.
Direct counting in $A,B,C,D$ gives each entry.
Every row except $(4,0)$ exceeds the flat cost $4x+24$ throughout
the stated interval. The $(4,0)$ tree has root children $G,\{4\},\{5\}$;
its $P$ and $Q$ costs are respectively
\[
 J_P(T_G)=\tfrac{49}{8}H+\tfrac{41}{2}S,\qquad
 J_Q(T_G)=\tfrac{31}{5}H+\tfrac{104}{5}S.
\]
The flat tree costs $4H+24S$ in both worlds.
Thus $T_G$ is uniquely preferred under $P$ and flat is uniquely preferred
under $Q$, with gaps $\Delta_P,\Delta_Q>0$.

For the minimax statement, larger trees cannot improve either action:
each contraction saves at least $m=2(H-S)$, and
$m-\Delta_P=(33H-44S)/8>0$ in the interval. Even a tree contracting to
$T_G$ is therefore componentwise dominated by flat. The only nondominated
actions are $T_G$ and flat. Choosing $T_G$ with probability $p$ has regrets
$(1-p)\Delta_P$ and $p\Delta_Q$; equating them proves the randomized
formula, attained at $p=\Delta_P/(\Delta_P+\Delta_Q)$.
\end{proof}

\begin{table}[t]\centering\footnotesize
\caption{Complete-epoch orbit certificate after contracting arities two and
three. Entries are total expected bytes divided by $S$, with $x=H/S$.
The flat comparator is $4x+24$.}\label{tab:strict-orbits}
\begin{tabular}{@{}crr@{}}\toprule
Child type & $P$ & $Q$\\\midrule
$(2,2)$ & $\frac{383}{48}x+\frac{335}{12}$ & $\frac{239}{30}x+\frac{418}{15}$\\
$(3,1)$ & $\frac{113}{16}x+\frac{97}{4}$ & $\frac{1131}{160}x+\frac{971}{40}$\\
$(3,2)$ & $8x+28$ & $\frac{639}{80}x+\frac{447}{16}$\\
$(4,0)$ & $\frac{49}{8}x+\frac{41}{2}$ & $\frac{31}{5}x+\frac{104}{5}$\\
$(4,1)$ & $\frac{113}{16}x+\frac{373}{16}$ & $\frac{71}{10}x+\frac{47}{2}$\\
\bottomrule\end{tabular}
\end{table}

At the artifact's $H=512,S=320$, $T_G$ costs 9,696 bytes under $P$
and $49,152/5=9,830.4$ under $Q$; flat costs 9,728.
The deterministic penalty is 32 bytes and randomized minimax regret is
$512/21\simeq24.381$ bytes, attained with group probability $5/21$.
An independent subset reference and separately written literal evaluator
agree on all 2,752 labeled unordered six-leaf trees, including uniqueness.
An exact affine certificate checks all two-node orbit classes across the
whole price interval. An integration test also fits the implemented full-wave
refiner on both exact mixtures and verifies its two different optima using
independent packet pricing; this is analytical conformance, not workload
performance evidence.

This is a constructed complete-epoch result, not an observed software
improvement. It allows three externally chosen refresh windows with
variable membership and width. The six-check CLI with its fixed eight-record
flush would instead produce one wave, so it does not instantiate this witness.
Other higher-order summaries, including wave-size distributions, may
distinguish these laws; the lower bound concerns the specified pairwise-only
information interface.


\section{Fixed-width padding and rank equivalence}
\subsection{Fixed-width batches retain a local information gap}
The previous theorem allows variable window widths. The following lift
connects the information limit to the evaluated eight-completion flush.
Its conclusion concerns two fixed feasible layouts; it does not assert
global optimality over all trees of the enlarged registry.

\begin{lemma}[Uniform equal-width padding]\label{lem:padding}
Let a random partition of $n$ original checks have $q$ blocks of size at
most $w$. Add $F=qw-n\ge2$ distinct filler checks, uniformly allocating
them to the block deficits. If $p_{ij}$ is the original probability that
$i,j$ share a block, then the same-block probabilities are
\begin{align}
 p_{if}&=\frac{w-1-\sum_{j\ne i}p_{ij}}{F},\label{eq:filler-one}\\
 p_{fg}&=\frac{qw(w-1)-(2w-2)n+2\sum_{i<j}p_{ij}}{F(F-1)}
 \quad(f\ne g).\label{eq:filler-two}
\end{align}
Thus identical original pairwise probabilities imply identical pairwise
probabilities on the padded registry.
\end{lemma}
\begin{proof}
If $i$ lies in a block of size $k$, its block receives $w-k$ fillers,
so a named filler joins it with probability $(w-k)/F$.
Taking expectations and using
$\E(k-1)=\sum_{j\ne i}p_{ij}$ gives~\eqref{eq:filler-one}.
For distinct fillers the conditional probability is
$\sum_{a=1}^q(w-k_a)(w-k_a-1)/[F(F-1)]$.
Expand the numerator and use $\sum_a k_a=n$ and
$\E\sum_a k_a(k_a-1)=2\sum_{i<j}p_{ij}$ to obtain
\eqref{eq:filler-two}. Original pair memberships are unchanged.
\end{proof}

\begin{proposition}[Fixed-eight candidate-class separation]\label{prop:fixed-eight}
For $b\ge6$ and the price interval of Theorem~\ref{thm:strict-pairwise},
there are two laws on complete 24-check execution sequences with identical
population bivariate completion-rank distributions, but opposite preferred
choices in a two-layout candidate class under consecutive batches of eight.
The two wrong-choice penalties remain $\Delta_P,\Delta_Q$; the class-relative
sharp randomized minimax regret is
$\Delta_P\Delta_Q/(\Delta_P+\Delta_Q)$.
\end{proposition}
\begin{proof}
Apply Lemma~\ref{lem:padding} to the six-check laws with $q=3,w=8,F=18$.
Their largest original block has size four. Original--filler probabilities
are $67/288$ for $G$ and $55/144$ for $O$; filler--filler probability is
$191/612$, in both worlds. Uniformly permute the three blocks and
independently shuffle the eight members of each. For two distinct checks
with same-block probability $p$, the probability of any specified pair of
distinct completion ranks is $p/[qw(w-1)]$ if the ranks lie in one batch,
and $(1-p)/[q(q-1)w^2]$ otherwise. Singleton ranks are uniform. Hence every
bivariate rank distribution agrees, while grouping consecutive ranks by
eight recovers the padded partition exactly.

Place the original six checks below one child of a common root and a fixed
feasible filler-only tree below its other child. The two candidate layouts
use respectively the flat six-leaf subtree and $T_G$ from
Theorem~\ref{thm:strict-pairwise}. The root and filler subtree contribute
the same cost to both choices within each world, including construction.
The original subtree's touched windows are unchanged. Subtracting the
common terms preserves $\Delta_P$ and $\Delta_Q$, and the same two-action
minimax calculation applies. Common costs need not agree between worlds.
\end{proof}

At the implemented prices, the wrong-choice penalties are 32 and 102.4
bytes and the class-relative randomized penalty is $512/21$ bytes.
The exact reference enumerates all 740 and 2,540 positive-mass padded
partitions in the smaller $w=4,F=6$ instance, independently checks the
pairwise formulas, and checks actual batching and packet costs on 600
complete 24-check sequences. Thirty tests cover these identities and
deliberate corruptions. Enumeration and representative sequence conformance support the
implementation; the padding and cancellation arguments establish the
general claim. Neither this construction nor its rank equivalence implies
that empirical pairwise-summary vectors have identical joint laws, or
that the two embedded layouts are globally optimal on 24 checks.


\section{Exact scoring and conditional adoption}
\subsection{Conditional adoption with lifecycle costs}\label{sec:adoption}
Selecting the smallest observed cost does not establish a net deployment
benefit. The following rule makes the additional assumptions explicit.
It specializes standard bounded-loss calibration and multiple-comparison
control~\cite{hoeffding,learntest}; it introduces no new concentration inequality.
Its HACT-specific step is to bound \emph{paired} differences after cancelling
packets with the same semantic coverage.

Fix an incumbent $L_0$ and $K\ge1$ training-fitted candidates $L_j$ before
calibration. All layouts use the same registry and packet contract.
An epoch partitions the $n$ checks into exactly $q$ nonempty waves of at
most $w$ checks, with $q\le n\le qw$. Let $h_e(G)$ count the waves that
intersect semantic coverage $G$. Aggregate the internal packet prices by
coverage and subtract the candidate's coefficients from the incumbent's:
\[
 a_{jG}=\sum_{v:G_{L_0}(v)=G}c(d_v)
          -\sum_{v:G_{L_j}(v)=G}c(d_v).
\]
Coefficients may have either sign. Shared packets cancel, regardless of
their display positions. The strict-epoch saving is exactly
\begin{align}\label{eq:paired-saving}
 d_j(e)&=\ell_e(L_0)-\ell_e(L_j)\nonumber\\
       &=a_{j0}+\sum_G a_{jG}h_e(G),\quad
 a_{j0}=\sum_Ga_{jG}.
\end{align}
For $s=|G|$, the support implies
\begin{equation}\label{eq:hit-envelope}
 l_G=\max\{\lceil s/w\rceil,q-(n-s)\}
       \le h_e(G)\le u_G=\min(s,q).
\end{equation}
The second lower bound follows because every untouched nonempty wave
requires a distinct check outside $G$. In particular $h_e(\mathcal R)=q$.
Define
\begin{align*}
 A_j&=a_{j0}+\sum_{a_{jG}>0}a_{jG}l_G
                    +\sum_{a_{jG}<0}a_{jG}u_G,\\
 B_j&=a_{j0}+\sum_{a_{jG}>0}a_{jG}u_G
                    +\sum_{a_{jG}<0}a_{jG}l_G.
\end{align*}
Set $R_j=B_j-A_j$. Thus $d_j(e)\in[A_j,B_j]$ on the entire declared support, rather than
merely the observed epochs. Bounds for different coverages need not be
jointly attainable, so this envelope is conservative.

Pairwise insufficiency does not require retaining every raw wave forever.
Once a catalogue is frozen, exact paired scoring needs only the accumulated
sum of each $d_j(e)$ and the number of epochs. More generally, the following
standard linear-identification argument characterizes population scoring
interfaces, and distinguishes exact score recovery from argmin recovery.

\begin{proposition}[Exact population score identification]\label{prop:summary-span}
On the finite epoch support $\Omega$, let an interface expose only
$\E_P\phi_1,\ldots,\E_P\phi_p$ for all probability laws $P$.
It determines every expected paired score $\E_Pd_j$, for every $P$, if and
only if each function $d_j$ belongs to
$\operatorname{span}\{1,\phi_1,\ldots,\phi_p\}$ on $\Omega$.
The smallest number of scalar expectation features for exact scoring is
the dimension of the span of the $d_j$ modulo constant functions.
\end{proposition}
\begin{proof}
Span membership supplies an affine recovery formula. Conversely, if
$d_j$ is outside the span, finite-dimensional orthogonality supplies a
signed vector $z$ annihilating $1$ and every $\phi_i$ but not $d_j$.
Its positive and negative parts have the same nonzero mass. Normalizing
them yields laws $P,Q$ with equal exposed summaries but different expected
paired scores. The dimension statement follows by taking a basis modulo
constants and applying the same argument to any lower-dimensional feature
span. This concerns exact scoring; different scores alone do not imply
different minimizing layouts, which requires the earlier separations.
\end{proof}

The implementation accumulates exact integer paired sums in one pass with
$O(K)$ score counters, in addition to the frozen coverage coefficients and
one epoch's masks. It retains no per-epoch loss rows on this path. Auditable
source/checker records and any required raw provenance remain separate;
score sufficiency grants no permission to discard acceptance evidence.

\begin{theorem}[Conditional expected lifecycle benefit]\label{thm:adoption}
Condition on training and the frozen catalogue, support, $0<\delta<1$,
TV radius $0\le\rho\le1$, horizon $h\ge1$ and nonnegative installation
costs $I_j$, all specified before calibration. Suppose the $m\ge1$
calibration epochs are IID from $P$ and independent of fitting. Let
$\bar d_j$ be their paired mean, and define
\begin{align}
 t_j&=R_j\sqrt{\frac{\log(K/\delta)}{2m}},\nonumber\\
 g_j&=\max\{A_j,\bar d_j-t_j-\rho R_j\}-I_j/h.
 \label{eq:adoption-lower}
\end{align}
With probability at least $1-\delta$ over calibration, simultaneously
for all candidates and every future law $Q$ on the same support with
$\operatorname{TV}(P,Q)\le\rho$,
\[
 h\,\E_Qd_j-I_j\ge h g_j.
\]
Adopting only a candidate with $g_j>0$, otherwise retaining $L_0$,
therefore certifies a conditional positive expected padded-byte benefit.
This does not certify realized savings, elapsed time, or task acceptance.
\end{theorem}
\begin{proof}
For each fixed candidate the one-sided Hoeffding inequality gives
$\Pr\{\E_Pd_j<\bar d_j-t_j\}\le\delta/K$ when $R_j>0$.
When $R_j=0$ the cost difference is constant and the conclusion holds
deterministically. A union bound makes these inequalities simultaneous;
candidate differences within an epoch need not be independent.

Use the convention $\operatorname{TV}(P,Q)=\sup_E|P(E)-Q(E)|$.
For $R_j>0$, normalize $(d_j-A_j)/R_j$ to $[0,1]$ and integrate
its superlevel-set probabilities. It follows that
$|\E_Qd_j-\E_Pd_j|\le\rho R_j$, uniformly over the TV ball.
For $R_j=0$ the difference is zero. Since $Q$ remains on the same
support, also $\E_Qd_j\ge A_j$. Combining the bounds and subtracting
$I_j/h$ proves the claim. The simultaneous event permits selection of
any qualifying candidate after calibration.
\end{proof}

The rule conditions on a \emph{justified} sampling design and shift radius;
neither is learned or verified from repeated observations by this theorem.
If deployment can leave the declared support, the guarantee does not apply.
Reusing calibration to refit layouts, choose the radius, or repeatedly
deploy new catalogues needs additional error accounting. The implementation
binds all plan parameters to a digest and requires external sampling/freeze
provenance; an assertion string alone is not evidence of independence or
chronology. Missing provenance or radius returns $\UNKNOWN$ and retains
the incumbent. Installation costs are bytes in the declared protocol;
compilation seconds and consumer utility remain separate admission tests.
An installation operation already charged inside the first epoch must not
also appear in $I_j$; the cost ledger must assign each operation once.


\section{Using the full-registry weights in the symmetric optimizer}
For declared batch sizes $w_1,\ldots,w_q$ summing to $n$, every node uses
\[
r_n(s)=1+\sum_a\left(1-\frac{\binom{n-s}{w_a}}{\binom n{w_a}}\right).
\]
The registry size remains $n$ inside every subtree state. Treating a subtree
as an independently completed registry defines a different law and can
change its optimum. Costs are additive across child subtrees; this
optimization decomposition does not assume probabilistic independence.

For example, let $n=4$, $b=2$, and $(w_1,w_2)=(2,2)$ at the implemented
prices $H=512,S=320$. A two-leaf internal node has weight $8/3$, while
three- and four-leaf internal nodes have weight $3$. A balanced binary
tree therefore costs $1152(3+2\cdot8/3)=9600$ expected bytes, including
construction. A binary chain costs $1152(3+3+8/3)=9984$. This is a
worked analytical example, not a measured software result.

The integer implementation scales rational weights by a common
denominator before comparing costs. Its arithmetic-operation bound is
$O(bn^2)$ after weights, with $O(bn)$ cost and split states; integer
bit complexity and reconstructing an output tree are additional accounting.
Independent enumeration checks 70 small cases and 74,245 feasible labeled
tree comparisons. These finite checks test the implementation; the induction
in the main article establishes the model-conditional general theorem.
The no-learning corollary additionally conditions on the future law given
training and selection. Unconditional exchangeability alone is insufficient.

\section{Evidence units and claim boundaries}
\begin{table}[htbp]
\centering\small
\caption{How the retained evidence supports the article. Re-encoding a
trace is not a new upstream execution.}
\begin{tabularx}{\linewidth}{@{}p{.25\linewidth}p{.25\linewidth}X@{}}
\toprule
Evidence & Unit or count & Supported interpretation\\\midrule
General information separation & 39 exact parameter cases &
Implementation conformance for a theorem proved for all admissible
fanouts and affine positive-slot prices.\\
Symmetric global reference & 70 cases; 74,245 tree comparisons &
Exact optimization under the declared completion law.\\
Public compiler cohort & Six requested projects; five admitted;
25 original admitted runs &
A convenience cohort of repeated suites, not independent source
changes or agent tasks. Packaging's resource exclusion remains.\\
Controlled dispatch shifts & 36 attempted processes; 30 admitted,
including six repaired runs &
Six initial VOID source-provenance outcomes are retained.
Artificial scheduling interventions do not establish natural drift.\\
Service observer ablation & 240 arms; 10,080 answer checks &
Local operation-time and answer conformance for fixed requests;
correlated timing blocks do not enlarge the project sample.\\
Retrospective symmetric scoring & The same 25 original and
30 admitted shift runs &
A stronger baseline comparison, with all original and file-group
losses retained; no additional checker execution.\\\bottomrule
\end{tabularx}
\end{table}

The symmetric layout beats frozen full-wave cap32 on all 15 item-shuffled
traces and loses on all 15 file-group traces. It costs 21.26--62.13\%
more on original repeated-order holdouts. With byte observers disabled,
the tree service takes 1.61--3.54 times the maintained-counter ledger's
operation time. No natural hierarchical consumer demand, general agent
speedup, token-billing reduction or physical-network benefit is inferred.

\section{Reproducing the frozen evidence}
The immutable scientific artifact is the public repository commit
\begin{center}\small
\texttt{ea874d49362552fd4c1369ef19fc307e6d040c7f}
\end{center}
available at
\begin{center}\small
\url{https://github.com/xuda1979/future-code/tree/ea874d49362552fd4c1369ef19fc307e6d040c7f/hact-paper/Adaptive_HACT_v5}.
\end{center}
The accompanying source archive contains the current article and supplement.
The observation transport in that artifact consists of 53 parts of a
31,138,565-byte lossless ZIP. Its SHA-256 is
\begin{center}\small\ttfamily
5b003c2566df152e77fef9e0af8f4caace4a446a9fa3bc5e7ce708feda65f35d
\end{center}
The index binds the original path, size and hash of each of 522 observations.
Hydration restores exact bytes, refuses mismatches and is offline.
The release inventory independently covers the restored observations as well
as their transport parts. Hashes provide byte integrity, not signatures.

With Python 3.12, use numpy 2.3.5, pytest 9.0.2 and pytest-asyncio 1.3.0.
From the artifact root:
\begin{verbatim}
python tools/hydrate_observations.py --require-index
python tools/verify_manifest.py
PYTEST_DISABLE_PLUGIN_AUTOLOAD=1 PYTHONPATH=. \
  python -m pytest -q -p pytest_asyncio.plugin
PYTHONPATH=. python -m experiments_v14.exchangeable_study --check
PYTHONPATH=. python -m audit.check_v14
PYTHONPATH=. python -m experiments_v14.report --check
\end{verbatim}
The reproduction guide lists all historical and current gates. Those checks
remain required: the primary workflow also rechecks v5/v9 arithmetic,
v10/v11 separation and exact references, v12 fixed batches and independent
follow-ups, v13 adoption and feature span, all manuscript projections,
the service oracle, measured query trace, and frozen source/synthetic
controls. The retained v12 audit report is
\path{.ai-loop/v12/reports/independent_followups.json}.
No missing observation should be repaired by changing its expected hash or
regenerating a measured record. A fresh computation goes to a new output path.

\section*{AI assistance}
OpenAI Codex (October 2026) assisted with this supplement's text, the
derivations and proof presentation, the worked example, and the associated
code and automated checks. These checks are distinct from human peer review.
The authors retain responsibility for the content.
\begin{thebibliography}{2}
\bibitem{hoeffding} W. Hoeffding. Probability Inequalities for Sums of
Bounded Random Variables. \emph{Journal of the American Statistical
Association}, 58(301):13--30, 1963.
\url{https://doi.org/10.1080/01621459.1963.10500830}.
\bibitem{learntest} A. N. Angelopoulos, S. Bates, E. J. Cand\`es,
M. I. Jordan, and L. Lei. Learn then Test: Calibrating Predictive
Algorithms to Achieve Risk Control. arXiv:2110.01052v5, 2022.
\url{https://arxiv.org/abs/2110.01052v5}.
\end{thebibliography}
\end{document}
